\exercisesection{§ 6 的习题}

在下面的习题中（Exerc. 3 除外），假设与记号均取自 §6。

\begin{enumerate}
\def\labelenumi{\arabic{enumi})}
\item
  假设 W 不可约。设 c 是 W 的一个 Coxeter 变换，\(\Gamma\) 是 W 中由 c 生成的子群，\(\Delta\) 是与 H 的某个元素正交的单位向量组成的集合。证明 \(\Gamma\) 在 \(\Delta\) 中有 l 个轨道，且每个轨道有 h 个元素。（按 Chap. VI, §1, no. 11 的 Prop. 33 的证明方式论证。）
\item
  假设 W 不可约。设 C 是相对于 W 的一个室，\((\mathrm{H}_{1},\ldots,\mathrm{H}_{l})\) 是它的墙，\(e_{i}\) 是与 \(H_{i}\) 正交的一个非零向量。假设 \(e_{1},\ldots,e_{r}\)（相应地，\(e_{r+1},\ldots,e_{l}\)）两两正交（cf.~no. 2）。对所有 \(u\in Z\)，定义 \(H_{u}\) 与 \(s_{u}\) 如下：\(H_{u}=H_{k}\)（若 \(u\equiv k\pmod{l}\)），且 \(s_{u}=s_{H_{u}}\)。
\end{enumerate}

\begin{enumerate}
\def\labelenumi{\alph{enumi})}
\item
  证明 \(\mathfrak{H}\) 的元素正是诸 \(s_1 s_2 \ldots s_{u-1} H_u\)，其中 \(u = 1, 2, \ldots, lh/2\)。
\item
  设 \(s' = s_1 \ldots s_r\)，\(s'' = s_{r+1} \ldots s_l\)，于是 \(c = s's''\) 是与有序室 C 相伴的 Coxeter 变换。设 \(w_0\) 是 W 中把 C 变到 -C 的元素（cf.~§4, Exerc. 2）。证明：若 \(h\) 为奇数，则
\end{enumerate}

\[
w _ {0} = \underbrace {s ^ {\prime} s ^ {\prime \prime} s ^ {\prime} \ldots s ^ {\prime \prime} s ^ {\prime}} _ {h \text {terms}} = \underbrace {s ^ {\prime \prime} s ^ {\prime} s ^ {\prime \prime} \ldots s ^ {\prime} s ^ {\prime \prime}} _ {h \text {terms}} = s ^ {\prime \prime} c ^ {(h - 1) / 2} = c ^ {(h - 1) / 2} s ^ {\prime}.
\]

\begin{enumerate}
\def\labelenumi{\alph{enumi})}
\setcounter{enumi}{2}
\tightlist
\item
  置 \(S = \{s_1, \ldots, s_l\}\)。偶 (W,S) 是一个 Coxeter 系统。若 \(w \in W\)，用 \(l_S(w)\) 表示 \(w\) 关于 S 的长度（Chap. IV, §1, no. 1）。证明
\end{enumerate}

\[
l _ {\mathrm{S}} (s ^ {\prime}) = r, l _ {\mathrm{S}} (s ^ {\prime \prime}) = l - r, l _ {\mathrm{S}} (c) = l.
\]

在 b) 的假设下，由此推出

\[
l _ {\mathrm{S}} (w _ {0}) \leqslant r + \frac {h - 1}{2} l \text { and } l _ {\mathrm{S}} (w _ {0}) \leqslant (l - r) + \frac {h - 1}{2} l.
\]

另一方面，证明 \(l_{S}(w_0) = \mathrm{Card}(\mathfrak{H}) = hl / 2\)（使用 Chap. IV, §1 的 Exerc. 22）。由此推出 \(r = l / 2\)，且 \((s_1, s_2, \ldots, s_{lh / 2})\) 是 \(w_0\) 的一个约化分解。

\(d)\) 证明：若 \(h\) 为偶数，则 \((s_1,\ldots ,s_{lh / 2})\) 是 \(w_{0} = c^{h / 2}\) 的一个约化分解。（方法同上。）

¶ 3) 设 K 是一个交换环，E 是一个以 \((e_1, \ldots, e_l)\) 为基的自由 K-模，\(f_1, \ldots, f_l\) 是 E 的对偶中的元素。置 \(a_{ij} = f_i(e_j)\)。若 \(1 \leqslant i \leqslant l\)，设 \(s_i\) 是伪反射 \(s_{e_i, f_i}\)（\S 2, Exerc. 1）。则

\[
s _ {i} (e _ {j}) = e _ {j} - a _ {i j} e _ {i}.
\]

置 \(c = s_1 \ldots s_l\)，\(z_i = c(e_i)\)。

\begin{enumerate}
\def\labelenumi{\alph{enumi})}
\tightlist
\item
  对 \(1 \leqslant i, k \leqslant l\)，置
\end{enumerate}

\[
y _ {i} ^ {k} = s _ {1} \dots s _ {k} (e _ {i}) \text { and } y _ {i} = s _ {1} \dots s _ {i - 1} (e _ {i}) = y _ {i} ^ {i - 1}.
\]

于是 \(y_{i}^{0} = e_{i}\)，\(y_{i}^{l} = z_{i}\)。

证明

\[
y _ {i} ^ {k - 1} - y _ {i} ^ {k} = a _ {k i} y _ {k}.
\]

由此推出公式

\[
e _ {i} = y _ {i} + \sum_ {k <   i} a _ {k i} y _ {k}
\]

\[
z _ {i} = y _ {i} - \sum_ {k \geqslant i} a _ {k i} y _ {k}.
\]

\begin{enumerate}
\def\labelenumi{\alph{enumi})}
\setcounter{enumi}{1}
\tightlist
\item
  设 C 是 \(c\) 关于基 \((e_i)\) 的矩阵。设 \(U = (u_{ij})\) 与 \(V = (v_{ij})\) 是由下式定义的矩阵
\end{enumerate}

\[
u _ {i j} = \left\{ \begin{array}{l l} a _ {i j} & \text { if } i <   j \\ 0 & \text { otherwise, } \end{array} \right. v _ {i j} = \left\{ \begin{array}{l l} 0 & \text { if } i <   j \\ a _ {i j} & \text { otherwise. } \end{array} \right.
\]

矩阵 \(\mathrm{I} + \mathrm{U}\) 可逆且行列式为 1。证明

\[
\mathrm{C} = (\mathrm{I} - \mathrm{V}) (\mathrm{I} + \mathrm{U}) ^ {- 1}.
\]

由此推出

\[
\det (\lambda I - C) = \det ((\lambda - 1) I + V + \lambda U),
\]

换句话说

\[
\det (\lambda \mathrm{I} - \mathrm{C}) = \left| \begin{array}{c c c c c} (\lambda - 1) + a _ {1 1} & \lambda a _ {1 2} & \lambda a _ {1 3} & \dots & \lambda a _ {1 l} \\ a _ {2 1} & \lambda (\lambda - 1) + a _ {2 2} & \lambda a _ {2 3} & \dots & \lambda a _ {2 l} \\ a _ {3 1} & a _ {3 2} & (\lambda - 1) + a _ {3 3} & \dots & \lambda a _ {3 l} \\ \vdots & \vdots & \vdots & \ddots & \vdots \\ a _ {l 1} & a _ {l 2} & a _ {l 3} & \dots & (\lambda - 1) + a _ {l l} \end{array} \right|.
\]

\begin{enumerate}
\def\labelenumi{\alph{enumi})}
\setcounter{enumi}{2}
\tightlist
\item
  设 \(\Gamma = (\mathrm{I},\mathrm{S})\) 是这样的图：其顶点集为 \(\mathrm{I} = [1,l]\)，其棱集 S 是 I 的二元子集 \(\{i,j\}\)（满足 \(a_{ij}\neq 0\) 或 \(a_{ji}\neq 0\)）组成的集合。若 \(\alpha = \{i,j\}\) 属于 S，用 \(\sigma_{\alpha}\) 表示 \(i\) 与 \(j\) 的对换（视为对称群 \(\mathrm{S_I}\) 的元素），并置 \(a_{\alpha} = -a_{ij}a_{ji}\)。
\end{enumerate}

设 \(\mathfrak{B}\) 是 \(S\) 中由顶点两两不交的棱组成的子集的集合。若 \(X \in \mathfrak{B}\)，用 \(C(X)\) 表示所有这样的 \(i \in I\)（它们不是 \(X\) 的任何棱的顶点）组成的集合，并置 \(\sigma_X = \prod_{\alpha \in X} \sigma_\alpha\)，\(a_X = \prod_{\alpha \in X} a_\alpha\)。

设 \(\sigma \in S_{I}\)，\(d_{\sigma}\) 是与 \(\sigma\) 对应的项（在 \((\lambda - 1)I + V + \lambda U\) 的行列式展开式中）。证明：若 \(\sigma\) 形如 \(\sigma_{X}\)，其中 \(X \in \mathfrak{B}\)，则

\[
d _ {\sigma} = a _ {X} \lambda^ {\operatorname{Card} (X)} \prod_ {i \in C (X)} (\lambda - 1 + a _ {i i}).
\]

现在假设 \(\Gamma\) 是一个森林（Chap. IV, Appendix, no. 3）。证明：若 \(\sigma \in S_l\) 不形如 \(\sigma_X\)（\(X \in \mathfrak{B}\)），则 \(d_{\sigma} = 0\)。由此推出公式

\[
\det (\lambda - c) = \sum_ {X \in \mathfrak {B}} a _ {X} \lambda^ {\operatorname{Card} (X)} \prod_ {i \in C (X)} (\lambda - 1 + a _ {i i}).
\]

\begin{enumerate}
\def\labelenumi{\alph{enumi})}
\setcounter{enumi}{3}
\tightlist
\item
  考虑多项式
\end{enumerate}

\[
P (X) = \left| \begin{array}{c c c c} X & a _ {1 2} & \ldots & a _ {1 l} \\ a _ {2 1} & X & \ldots & a _ {2 l} \\ \vdots & \vdots & \ddots & \vdots \\ a _ {l 1} & a _ {l 2} & \ldots & X \end{array} \right|.
\]

在 \(c\) ) 的假设之外，再加设 \(a_{ii} = 1\)（对所有 \(i\)）。证明在此情形下

\[
\det (\lambda^ {2} - c) = \lambda^ {l} P (\lambda + \lambda^ {- 1}).
\]

\begin{enumerate}
\def\labelenumi{\arabic{enumi})}
\setcounter{enumi}{3}
\tightlist
\item
  在 no. 1 的假设与记号下，用 \(n_{ij}\) 表示 \(s_{\mathrm{H}_i}s_{\mathrm{H}_j}\) 的阶，并置 \(a_{ij} = -2\cos \frac{\pi}{n_{ij}}\)。于是 \(a_{ii} = 2\)，且 \(a_{ij} \leqslant 0\)（若 \(i \neq j\)），cf.~§3。利用前一习题的方法证明 \(^{12}\)：
\end{enumerate}

\[
\left| \begin{array}{c c c c} X & a _ {1 2} & \ldots & a _ {1 l} \\ a _ {2 1} & X & \ldots & a _ {2 l} \\ \vdots & \vdots & \ddots & \vdots \\ a _ {l 1} & a _ {l 2} & \ldots & X \end{array} \right| = \prod_ {i = 1} ^ {l} (X - 2 \cos (\pi m _ {i} / h)),
\]

其中 \(m_{1},\ldots ,m_{l}\) 是 W 的指数。

